\chapter{Návrh rozhraní}
\label{chap:design}

V předchozí kapitole jsme popsali možnosti a omezení původních rozhraní
generátorů \texttt{geng} a \texttt{plantri}. V této kapitole představíme návrh
knihovny PlanG, která oba generátory zpřístupňuje prostřednictvím společného
rozhraní v jazyce C++.

Zaměříme se na rozdělení společných a specifických částí rozhraní, registraci
callbacků, zpřístupnění grafů bez jejich kopírování a integraci s Boost Graph
Library. Technické propojení s původními C programy popíšeme v kapitole
o implementaci.

\section{Cíle návrhu}
\label{sec:design-goals}

Původní generátory použijeme jako backendy a nad nimi nabídneme společné C++
rozhraní. Uživatel si zvolí požadovaný backend, nastaví parametry generování pomocí
typově bezpečných metod a připojí vlastní funkce pro filtrování, prořezávání
nebo zpracování výsledků. Společné operace mají být pro oba generátory stejné,
zatímco nastavení závislá na konkrétním generátoru zůstávají součástí jeho
backendu.

Důležitým požadavkem je zachování výkonu původních programů. Rozhraní proto
nemá automaticky převádět každý vygenerovaný graf do nové datové struktury.
Místo toho graf zpřístupní pomocí view nad interními daty generátoru. Uživatel
tak může graf zpracovat přímo v callbacku bez vytvoření jeho úplné kopie.

Návrh má také umožnit použití běžných grafových algoritmů a případné doplnění
dalších generátorů. Rozhraní proto odděluje společné řízení generování od
vlastností konkrétního backendu a přizpůsobuje grafová views vybraným
konceptům Boost Graph Library.

\section{Jednotné rozhraní pro více generátorů}
\label{sec:unified-interface}


Základním vstupním bodem knihovny je šablona
\texttt{Generator\textless Backend\textgreater}. Typ předaný jako
\texttt{Backend} určuje, který původní generátor se použije a jaká specifická
nastavení budou dostupná. Pro obecné grafy použijeme typ
\texttt{Generator\textless geng::Backend\textgreater}, zatímco pro rovinné
grafy typ \texttt{Generator\textless plantri::Backend\textgreater}.

Šablona \texttt{Generator} poskytuje společné operace, zejména spuštění
generování metodou \texttt{run} a registraci callbacků. Zvolený backend
doplňuje metody pro konfiguraci vlastností podporovaných konkrétním
generátorem. Uživatel tak pracuje s oběma generátory stejným způsobem, ale
zároveň má přístup k jejich specifickým možnostem.

Toto uspořádání využívá návrh založený na politikách (policy-based design).
Konkrétní backend je zvolen jako parametr šablony a poskytuje třídě 
\texttt{Generator} chování závislé na daném generátoru. Volba backendu tak 
probíhá při překladu a nevyžaduje virtuální funkce.

Backend musí poskytovat typ grafového view, vlastní konfiguraci a operace
potřebné ke spuštění generátoru. Díky tomuto rozdělení lze knihovnu rozšířit
o další generátor vytvořením nového backendu, aniž by bylo nutné měnit způsob,
jakým uživatel spouští generování nebo registruje callbacky.

\section{Callbacky pro filtrování a zpracování grafů}
\label{sec:design-callbacks}

V rozhraní PlanG rozlišujeme tři typy callbacků, které odpovídají různým
fázím generování. Registrujeme je jako C++ funkční objekty před spuštěním
metody \texttt{run}. Můžeme přitom použít běžnou funkci, lambda výraz nebo
objekt s operátorem volání.

Pomocí metody \texttt{setPrune} filtrujeme výsledné grafy. Callback pro každý
dokončený graf vrátí hodnotu \texttt{KEEP}, pokud jej chceme ponechat, nebo
\texttt{PRUNE}, pokud jej chceme odmítnout.

Metodu \texttt{setPreprune} používáme pro prořezávání rozpracovaných větví.
Callback rozhoduje, zda má generátor pokračovat v aktuální větvi. Protože se
může volat mnohem častěji než \texttt{setPrune}, měl by prováděný test být
rychlý a nesmí odmítnout větev, která by mohla vést k přípustnému výsledku.

Pomocí metody \texttt{setOutproc} zpracováváme přijaté výsledné grafy. Můžeme
ji využít pro vlastní výstup, sběr statistik nebo spuštění dalšího grafového
algoritmu.

Do callbacků předáváme objekt \texttt{GraphView}, který zpřístupňuje aktuální
graf bez vytvoření jeho úplné kopie. S grafem tak pracujeme stejným způsobem
nezávisle na zvoleném backendu.

\section{\texttt{View} na graf bez kopírování dat}
\label{sec:graph-view}



Při zpracování velkého počtu grafů nechceme každý výsledek převádět do nové
vlastněné datové struktury. Pro každý backend proto definujeme samostatný typ
\texttt{GraphView}, který zpřístupňuje graf uložený uvnitř původního
generátoru.

\texttt{geng::GraphView} uchovává ukazatel na bitovou reprezentaci nauty,
zatímco \texttt{plantri::GraphView} pracuje s poli stupňů a orientovanými
hranami rovinného vložení. Vytvořením objektu \texttt{GraphView} kopírujeme
pouze několik ukazatelů a základních údajů o grafu, nikoli jeho vrcholy a
hrany.

Uvnitř callbacku pracujeme s grafem zpravidla prostřednictvím poskytovaných
free funkcí, například \texttt{vertices}, \texttt{edges},
\texttt{adjacent\_vertices} nebo \texttt{out\_edges}. Tyto funkce tvoří
rozhraní požadované vybranými koncepty Boost Graph Library. Můžeme je používat
přímo nebo objekt \texttt{GraphView} předat kompatibilnímu algoritmu z této
knihovny.



\section{Integrace s Boost Graph Library}
\label{sec:boost-integration}

Boost Graph Library je založena na grafových konceptech, které určují, jaké
typy a operace musí graf poskytovat, aby nad ním bylo možné spouštět obecné
grafové algoritmy~\cite{boostGraphConcepts}. Graf přitom nemusí být instancí
konkrétní třídy z knihovny Boost, pokud splňuje požadované rozhraní.

Pro oba typy \texttt{GraphView} definujeme typy vrcholů a hran, iterátory
a podporované způsoby procházení grafu. Společně s free funkcemi z předchozí
sekce tak oba typy splňují vybrané koncepty Boost Graph Library.

Díky této integraci můžeme v callbacku používat například algoritmy pro
barvení vrcholů nebo testování rovinnosti. Konkrétní implementaci potřebných
typů, iterátorů a traits popíšeme v kapitole o implementaci.

\section{Podpora bohatších struktur}
\label{sec:richer-structures}

Veřejné rozhraní backendu \texttt{geng} podporuje generování grafů s barevnými
třídami vrcholů. Uživatel může nastavit počet barev, přesné velikosti tříd
nebo jejich dolní a horní meze pomocí metod \texttt{setColors},
\texttt{setColorClassSizes} a \texttt{setColorClassBounds}.

Zakotvené vrcholy zadáváme metodou \texttt{setRootedVertices}. Každý
zakotvený vrchol tvoří samostatnou jednoprvkovou barevnou třídu, kterou musí
izomorfismus zachovat.

Přiřazené barvy zpřístupňujeme také prostřednictvím
\texttt{geng::GraphView}, aby s nimi mohly pracovat uživatelské callbacky.
Způsob generování barev a jejich zahrnutí do kanonických testů popíšeme
v kapitole o implementaci.
